% ============================================================
%  P3 — Cross-Material Transfer of GLAD Column-Tilt and
%       Porosity Models: A Rigorous Negative Result and Its
%       Statistical Boundary
%
%  Template: AIP Publishing (revtex4-1, aip class option)
%  Target journal: AIP Advances (companion submission to P1/P2,
%          same publisher family)
%
%  Key results (all numbers traced to 02_PINN_NIF/P3_ML/*.md,
%  none recomputed for this draft):
%    - Pooled LOMO-CV cross-material model for column-tilt beta(alpha),
%      n=178 rows / 37 materials: gradient-boosted trees +49.2% RMSE
%      vs. tangent-rule physics baseline, permutation p=0.22 (n_perm=100).
%      RMSE bar cleared, significance bar NOT cleared.
%    - Porosity(alpha), n=91/19 (post-audit corpus): random forest
%      +19.2% RMSE vs naive-mean baseline, p=0.36. NOT cleared.
%    - refractive_index(alpha), n=78/13: GP +16.2%, p=0.54. NOT cleared.
%    - band_gap Eg(alpha), n=67/10: no model beats naive-mean baseline.
%      NOT cleared (fails RMSE bar too).
%    - Material-class stratification (elemental metals/oxides-fluorides/
%      chalcogenides): no subgroup clears p<0.05; elemental metals
%      closest (p=0.20 at n=14).
%    - Targeted elemental-metals corpus growth (cross-source and
%      same-source densification): p ranges 0.20-0.36 across every
%      variant tried; no lever reliably improves it.
%    - Pairwise-ranking reformulation: beta accuracy 0.746 (vs 0.544
%      baseline), p=0.27; porosity accuracy 0.725 (vs 0.602), p=0.28.
%      Both clear the accuracy bar, neither clears significance.
%    - Bayesian hierarchical partial-pooling model (n=213/40): atomic
%      mass's effect on the intercept excludes zero (94% HDI
%      [-7.24,-0.34]); Bayesian R^2 median 0.28 (intercept layer),
%      0.20 (slope layer); P(R^2>0.1) = 92%/80%. Genuine complementary
%      positive finding, not in conflict with the frequentist verdict.
%
%  No figures: all results are reported as tables (this simplifies
%  the AIP conversion relative to P1/P2, which each carry a flat set
%  of figure PDFs).
%
%  Converted 2026-08-20 from the plain-article-class source
%  (05_THESIS_PAPER_ASSETS/article_draft/P3_cross_material_generalization_GLAD.tex)
%  to AIP's revtex4-1 template, at the user's request, targeting AIP
%  Advances (companion submission to P1/P2, same publisher family).
%  Content, tables, and references are unchanged; only the document
%  class, title-block macros, and back matter were converted.
%  Bibliography style switched to aipnum4-1.bst (numbered, AIP house
%  style); \citep/\citet unchanged (revtex4-1 loads natbib internally).
%  Author: Oussama Ben Haj Ali (LPMS / ENIT / UTM)
% ============================================================
\documentclass[aip,amsmath,amssymb,reprint]{revtex4-1}

\usepackage{graphicx}
\usepackage{booktabs}
\usepackage{siunitx}
\usepackage{xcolor}
\usepackage{nicefrac}

\renewcommand{\ang}{\ensuremath{\alpha}} % bare-symbol form used throughout,
                                          % as in the companion P1/P2 manuscripts
\newcommand{\lomo}{\textsc{lomo-cv}}
\newcommand{\beto}{\ensuremath{\beta(\alpha)}}
\newcommand{\poro}{\ensuremath{\phi(\alpha)}}

% ============================================================
\begin{document}

\title{Cross-Material Transfer of GLAD Column-Tilt and Porosity
  Models: A Rigorous Negative Result and Its Statistical
  Boundary}

\author{Oussama Ben Haj Ali}
\email{oussama.benhadjali@fst.utm.tn}
\affiliation{Laboratoire de Photovolta\"{i}que et Mat\'{e}riaux
  Semi-conducteurs (LPMS), \'{E}cole Nationale d'Ing\'{e}nieurs de Tunis (ENIT),
  Universit\'{e} de Tunis El Manar (UTM), Tunis, Tunisia}

\author{Ferid Chaffar Akkari}
\affiliation{Laboratoire de Photovolta\"{i}que et Mat\'{e}riaux
  Semi-conducteurs (LPMS), \'{E}cole Nationale d'Ing\'{e}nieurs de Tunis (ENIT),
  Universit\'{e} de Tunis El Manar (UTM), Tunis, Tunisia}

\author{Bruno Gallas}
\affiliation{Institut des NanoSciences de Paris (INSP), CNRS,
  Sorbonne Universit\'{e}, Paris, France}

\author{Nicolas Martin}
\affiliation{FEMTO-ST Institute, CNRS, Universit\'{e} de Franche-Comt\'{e},
  Besan\c{c}on, France}

\date{\today}

\begin{abstract}
Glancing-angle deposition (GLAD) produces tilted, porous columnar thin
films whose column-tilt angle $\beta$ and porosity $\phi$ are the two
morphological quantities most directly linked to a film's optical and
electrical anisotropy. Existing predictive relations for both
quantities are either purely geometric (the tangent rule for $\beta$)
or material-specific empirical fits; no published model uses material
identity itself -- through physically motivated descriptors such as
atomic mass, melting point, or bulk density -- to predict how a
\emph{new} material, never seen during fitting, will behave under
GLAD. We assembled a literature-derived corpus of experimentally
measured $\beta(\alpha)$ and porosity($\alpha$) values spanning up to
40 distinct materials (elemental metals, oxides, fluorides,
chalcogenides), fitted four candidate models (linear regression,
Gaussian process regression, gradient-boosted trees, random forest)
under leave-one-material-out cross-validation (\lomo{}), and
pre-registered a strict acceptance bar before inspecting any result:
$\geq 15\%$ \lomo{} RMSE improvement over a physics/naive baseline
\emph{and} a material-level permutation-test $p<0.05$. Every learned
model comfortably clears the RMSE bar for $\beta(\alpha)$ (up to
$+54.4\%$ depending on corpus snapshot) and often for porosity($\alpha$)
(up to $+28.5\%$), but no configuration -- the pooled fit, three
pre-registered material-class subgroups, two independent corpus-growth
interventions targeting the most promising subgroup, or a pairwise
ranking reformulation of the task -- ever clears the permutation-test
bar (best $p=0.09$ for $\beta(\alpha)$ at an earlier, smaller corpus
snapshot; $p=0.20$--$0.36$ across every variant tested at the current
snapshot). We report this as an honest, thoroughly instrumented
negative result: the barrier is specifically material-level sample
size, not model choice, feature choice, task formulation, or
extraction quality, all four of which were tested and ruled out as
explanations. As a complementary, deliberately different analysis, a
pre-registered Bayesian hierarchical partial-pooling model fitted to
the same $\beta(\alpha)$ corpus ($n=213$ rows, $40$ materials, all MCMC
diagnostics clean) finds that atomic mass has a credible non-zero
effect on a material's mean offset from the tangent rule (94\% highest
density interval excluding zero) and that Bayesian $R^2$ posteriors
place 80--92\% probability on material descriptors explaining more
than 10\% of between-material variance. This does not overturn the
frequentist non-significance verdict -- it explains it: a
real-but-modest effect of the size the Bayesian model finds is exactly
what a permutation test at $n=14$--$40$ materials is underpowered to
detect at $p<0.05$. We present this paper as a resource for the
GLAD/materials-informatics community: a fully pre-registered,
diagnostically exhaustive account of where the statistical-power
boundary for this class of cross-material generalization problem
currently sits, together with the specific corpus-growth target
(professional literature-database access to 30+ elemental metals,
particularly platinum-group and rare-earth species) that would most
directly move it.
\end{abstract}

\maketitle

% ============================================================
\section{Introduction}
\label{sec:intro}

Glancing-angle deposition (GLAD) is a physical vapor deposition
technique in which the substrate is tilted at a large angle
$\alpha$ (measured from the surface normal) relative to the incoming
vapor flux, so that geometric self-shadowing among growing nuclei
produces isolated, tilted columnar structures rather than a dense
film \citep{robbie1997,hawkeye2007,hawkeye2014}. The resulting films
combine two properties of independent technological interest: a
column-tilt angle $\beta$, generally smaller than the deposition angle
$\alpha$ because incoming flux is captured preferentially by the
upper, unshadowed part of a growing column, and an internal porosity
$\phi$ set by how completely that self-shadowing excludes flux from
the space between columns. Both properties drive the strong optical
and electrical anisotropy that makes GLAD films attractive for
polarization optics, sensing, and photovoltaic light management
\citep{brett2008science,lakhtakia2005stf,messier2000origin}.

The relationship between $\alpha$ and $\beta$ has a long-standing
geometric approximation, the tangent rule
\citep{tait1993modelling,thornton1974}, which treats column growth as
a purely ballistic shadowing process independent of the depositing
material's identity. In practice, measured $\beta(\alpha)$ curves for
different materials deposited under nominally comparable conditions
show real, reproducible scatter around the tangent-rule prediction
\citep{chaffar2018,chaffar2019,tounsi2015,benamor2026,%
benNacer2020cu2oGlad} -- a scatter plausibly attributable to
material-dependent surface mobility, sticking coefficient, and
crystal structure, none of which the tangent rule encodes. Porosity
has no comparably established closed-form geometric rule at all;
published porosity($\alpha$) values are almost always reported and
modeled per-material. This project's own ballistic GLAD simulator
(companion manuscript, \citealp{companionP1}) and density-calibrated
neural-field surrogate (companion manuscript, \citealp{companionP2})
both operate at fixed, explicitly specified material parameters; a
model that could predict $\beta(\alpha)$ or $\phi(\alpha)$ for a
material \emph{not} already in the simulator's parameter database,
using only readily available descriptors of that material, would be
directly useful for screening new material choices before running any
simulation or deposition campaign.

This motivates a natural question: can material descriptors
(atomic mass, melting point, bulk density, crystal structure,
deposition method) predict a material's $\beta(\alpha)$ or
$\phi(\alpha)$ curve well enough, and reliably enough across held-out
materials, to be useful for this kind of screening? We built a
literature-derived corpus spanning up to 40 distinct materials across
elemental metals, oxides, fluorides, and chalcogenides, and answered
this question with a pre-registered, leakage-controlled
leave-one-material-out (\lomo{}) evaluation protocol, reported here in
full together with the extensive, honestly documented attempts made to
recover a positive result once the pooled result fell short. We regard
the resulting negative-to-mixed finding as a genuine contribution in
its own right, not a shortfall to be minimized: it identifies exactly
where, quantitatively, the statistical-power boundary for this class
of materials-informatics problem currently sits, rules out four
plausible alternative explanations (wrong features, wrong task
formulation, class heterogeneity, extraction quality) for the
boundary, and provides a Bayesian complementary analysis that
characterizes the size of the effect the frequentist test cannot yet
resolve. This is offered explicitly as a resource for future
corpus-growth efforts in this specific sub-field, in the spirit that a
well-instrumented negative result that saves the next researcher from
repeating the same five failed levers is worth reporting on its own
merits.

The remainder of this paper is organized as follows.
Section~\ref{sec:methods} describes the corpus, the pre-registered
evaluation protocol, the four candidate models, the material-class
stratification, the pairwise-ranking reformulation, and the Bayesian
hierarchical model specification. Section~\ref{sec:results} reports
the pooled results for all four target observables, the
class-stratified results, the targeted corpus-growth push within the
most promising class, the ranking reformulation, and the Bayesian
hierarchical analysis. Section~\ref{sec:discussion} interprets these
results together, argues for reading them as a genuine
statistical-power finding rather than a failure of modeling
approach, and identifies the concrete corpus-growth target that would
most directly resolve the open question. Section~\ref{sec:conclusion}
summarizes.

% ============================================================
\section{Methods}
\label{sec:methods}

\subsection{Corpus and target observables}
\label{ssec:corpus}

Four target observables were evaluated, all as a function of
deposition angle $\alpha$: column-tilt angle $\beta(\alpha)$,
porosity $\phi(\alpha)$, refractive index $n(\alpha)$, and optical
band gap $E_g(\alpha)$. Rows were drawn from the project's ongoing
literature-absorption effort (peer-reviewed papers and, in a small
number of cases, PhD theses, reporting angle-resolved measurements of
one or more of these quantities for a named material) and assembled
into per-observable feature tables by an automated pipeline
(\texttt{build\_feature\_table.py}) that attaches five candidate
material descriptors to every row: atomic mass, melting point, bulk
density, crystal structure (categorical), and deposition method
(categorical: e-beam evaporation, thermal evaporation, DC/RF
sputtering, etc.). Rows lacking a numeric target value, an unresolved
deposition angle, or a required descriptor for the base material were
excluded (26/228 excluded rows for $\beta$, 11/114 for porosity, at the
snapshot reported in Table~\ref{tab:coverage}); exclusion reasons are
logged per-row and are not adjusted after seeing model results. Corpus
size grew substantially over the course of this investigation as
literature absorption continued in parallel with model evaluation;
every result below states the exact $(n_{\mathrm{rows}},
n_{\mathrm{materials}})$ snapshot it was computed on, since corpus
snapshots differ across sections of this paper by design -- later
sections report tests run against a larger corpus than earlier
sections, reflecting real, ongoing corpus growth rather than
selective reporting.

\begin{table}[htbp]
  \centering
  \caption{Feature-table coverage at the most complete raw-table
    snapshot produced during this investigation (2026-08-17), shown
    here to illustrate typical exclusion reasons and per-observable
    coverage. Because literature absorption continued throughout the
    investigation, this raw snapshot is not identical to any single
    downstream analysis's corpus: the pooled baseline
    (Section~\ref{ssec:pooled}), the earlier corpus-growth campaign
    (Table~\ref{tab:corpus_growth}), the class-stratification and
    ranking analyses (Sections~\ref{ssec:stratification}--\ref{ssec:ranking}),
    and the Bayesian hierarchical model (Section~\ref{ssec:bayes}) were
    each run at their own specific, smaller snapshot along this same
    growth trajectory (after an analysis-specific descriptor-completeness
    filter), and each result below states its own exact
    $(n_{\mathrm{rows}}, n_{\mathrm{materials}})$.}
  \label{tab:coverage}
  \begin{tabular}{lrr}
    \toprule
    Target observable & Rows & Distinct materials \\
    \midrule
    Column-tilt angle $\beta(\alpha)$ & 228 (202 usable) & 43 \\
    Porosity $\phi(\alpha)$ & 114 (103 usable) & 23 \\
    Refractive index $n(\alpha)$ & 105 (102 usable) & 16 \\
    Band gap $E_g(\alpha)$ & 78 (72 usable) & 12 \\
    \bottomrule
  \end{tabular}
\end{table}

\subsection{Pre-registered evaluation criterion}
\label{ssec:preregistration}

Before fitting any model to a given corpus snapshot, a single
acceptance criterion was fixed for each target observable,
independent of the outcome: a candidate model was considered to have
cleared the bar only if it achieved (i) a leave-one-material-out
cross-validated RMSE at least 15\% lower than the corresponding
baseline (the tangent rule for $\beta(\alpha)$; the \lomo{} train-mean
for porosity, refractive index, and $E_g$, since no established
closed-form geometric rule exists for these), and (ii) a
material-level permutation-test $p$-value below 0.05. Both conditions
were fixed prior to inspecting either result and applied identically
across all four observables and every corpus snapshot reported below.
Results are reported as obtained, including every negative and
partially negative outcome, without post-hoc adjustment of the
criterion.

\subsection{Leave-one-material-out cross-validation}
\label{ssec:lomo}

All model comparisons use leave-one-material-out (\lomo{}) rather than
leave-one-row-out cross-validation: rows from the same material
(often multiple angles from a single reproducibility series or
multiple points digitized from one figure) are not statistically
independent, and evaluating a model on a held-out \emph{row} whose
material was seen during training would leak information the model
could exploit without any genuine cross-material generalization
ability. Under \lomo{}, one material's entire row set is held out as
the test fold on each iteration, and the model is refit on the
remaining materials; this directly measures whether descriptor-based
generalization to a truly unseen material is possible, which is the
scientific question this paper asks.

\subsection{Candidate models}
\label{ssec:models}

Four models were compared under the identical \lomo{} protocol, plus
the relevant baseline for each observable:
\begin{itemize}
  \item \textbf{Physics/naive baseline}: the tangent rule (a
    zero-free-parameter geometric relation) for $\beta(\alpha)$; the
    \lomo{} train-mean for porosity, refractive index, and $E_g$
    (no established closed-form rule for these).
  \item \textbf{Linear regression} on the standardized descriptor set.
  \item \textbf{Gaussian process regression} (Mat\'{e}rn-5/2 kernel plus
    a white-noise term; \texttt{scikit-learn}
    \texttt{GaussianProcessRegressor}) \citep{rasmussen2006gp,%
    pedregosa2011sklearn} -- designated the primary candidate at the
    project's earliest evaluation round, before the gradient-boosted
    and random-forest comparators were added.
  \item \textbf{Gradient-boosted trees}
    \citep{friedman2001gbm,pedregosa2011sklearn}.
  \item \textbf{Random forest} \citep{breiman2001randomforest,%
    pedregosa2011sklearn}.
\end{itemize}
All descriptor-based models share the same feature pipeline: numeric
descriptors are standardized (zero mean, unit variance, fit on the
training fold only) and categorical descriptors (crystal structure,
deposition method) are one-hot encoded, both wrapped in a
\texttt{scikit-learn} \texttt{ColumnTransformer}/\texttt{Pipeline} so
that no fold-specific statistic is computed on held-out data.

\subsection{Material-level permutation test}
\label{ssec:permutation}

Statistical significance of a model's \lomo{} RMSE improvement over
its baseline is assessed with a material-level permutation test
\citep{ojala2010permutation}: the
material-block labels attached to each row (not individual rows) are
shuffled $n_{\mathrm{perm}}=100$ times, the same \lomo{} model-fitting
and RMSE-improvement calculation is repeated on each shuffled
assignment, and $p$ is the fraction of shuffles whose RMSE improvement
equals or exceeds the improvement observed on the real, unshuffled
data. Shuffling at the material-block level, rather than the row
level, is essential: row-level shuffling would break the same
within-material correlation structure that motivates \lomo{} in the
first place and would understate the true null distribution's spread,
inflating apparent significance. An early corpus snapshot ($n=73$
rows, 23 materials) returned $p=0.020$ for $\beta(\alpha)$, initially
reported as a cleared result; this did not survive further, real
corpus growth (Section~\ref{ssec:pooled}), which is itself part of
this paper's evidence that a single small-corpus permutation-test
result should not be treated as final without a subsequent replication
check as the corpus grows.

\subsection{Material-class stratification}
\label{ssec:stratification_method}

To test whether the pooled model's failure to clear significance
reflects genuine cross-class physical heterogeneity being averaged
away -- elemental metals, ionic/covalent oxides and fluorides, and
covalent chalcogenides plausibly respond differently to ballistic
shadowing -- a grouping by crystal-structure/bonding chemistry was
pre-registered and frozen \emph{before} any per-class fit was run:
elemental metals (fcc/bcc/hcp pure elements), rock-salt nitrides
(grouped with AlN by anion chemistry), oxides/fluorides
(ionic/covalent, ZnO included by anion chemistry despite its wurtzite
structure), and chalcogenides/sulfides. A minimum of five materials
per subgroup was pre-registered as a floor for testing; subgroups
below this floor are reported as untested, not as null results. Each
qualifying subgroup was fit with the identical four-model comparison,
\lomo{}-CV, and material-level permutation protocol used for the
pooled model.

\subsection{Pairwise ranking reformulation}
\label{ssec:ranking_method}

As an orthogonal test of whether the pre-registered exact-value
regression criterion is testing the wrong task, a pairwise-ranking
reformulation was pre-registered
(\texttt{RANKING\_REFORMULATION\_PREREGISTRATION\_20260817.md}): given
two materials A and B measured at angles within $\pm 5^{\circ}$ of
each other, predict whether A's target value exceeds B's, rather than
predicting either value exactly. This is scientifically useful in its
own right (materials-selection screening does not always require an
exact predicted value, only a reliable ranking) and might expose
cross-material signal that exact-value regression cannot. The same
\lomo{} and material-level permutation protocol was applied to the
ranking task ($n_{\mathrm{perm}}=100$), against a separately
pre-registered bar: pooled \lomo{} pairwise accuracy $\geq 0.60$
\emph{and} permutation $p<0.05$. Candidate models were a zero-parameter
tangent-rule baseline, an angle-difference-only logistic baseline (to
isolate any signal carried by $\alpha$ alone, independent of material
identity), a full-descriptor linear (logistic) ranker, and a
full-descriptor gradient-boosted ranker.

\subsection{Bayesian hierarchical partial-pooling model}
\label{ssec:bayes_method}

As a final, deliberately different analysis, a Bayesian hierarchical
(partial-pooling) model was pre-registered
(\texttt{BAYESIAN\_HIERARCHICAL\_PREREGISTRATION\_20260817.md}, fixed
before fitting) and fit to the $\beta(\alpha)$ corpus using
\texttt{PyMC} 5.28.5 with \texttt{ArviZ} 0.23.4 for diagnostics
\citep{salvatier2016pymc3}, sampled with the No-U-Turn Sampler (NUTS,
CPU only). This model asks a different question from the frequentist
tests above: not ``does a flexible model out-predict the physics
baseline on a held-out material by a margin unlikely under
label-shuffling?'', but ``how much of the material-to-material spread
in $\beta(\alpha)$ is attributable, in a regularized partial-pooling
regression, to physically motivated descriptors, and with what
uncertainty?'' The model places a per-material random intercept
$a_i$ and slope $b_i$ on the tangent-rule residual (i.e., a
material-specific affine correction to the physics baseline as a
function of $\alpha$), regressed at the population level on three
standardized descriptors (atomic mass, melting point, bulk density);
crystal structure was excluded from the group-level regression for
identifiability reasons at 19 levels across 43 materials. Convergence
was assessed with the rank-normalized $\hat{R}$ diagnostic
\citep{vehtari2021rhat} and effective sample size. The primary
deliverable is the Bayesian $R^2$
\citep{gelman2019r2} -- computed per posterior draw as
$\mathrm{Var}(\mathrm{fitted})/(\mathrm{Var}(\mathrm{fitted}) +
\mathrm{residual\ variance})$, separately for the intercept layer
($R^2_a$) and the slope layer ($R^2_b$) -- reported as an in-sample
quantity (how much structure the fitted model captures), not a
cross-validated one (how well it would predict a new material).

% ============================================================
\section{Results}
\label{sec:results}

\subsection{Pooled cross-material models: RMSE bar cleared, significance bar not}
\label{ssec:pooled}

Table~\ref{tab:pooled_all} summarizes the best-model \lomo{} result
for all four target observables at the post-comprehensive-audit corpus
snapshot ($\beta(\alpha)$: $n=178$ rows/37 materials;
porosity($\alpha$): $n=91/19$; refractive index: $n=78/13$; $E_g$:
$n=67/10$). For $\beta(\alpha)$, the tangent-rule baseline achieves a
\lomo{} RMSE of 27.136$^{\circ}$; every learned model outperforms it,
with gradient-boosted trees the best performer at 13.782$^{\circ}$
($+49.2\%$). This comfortably clears the pre-registered 15\% RMSE bar.
The material-level permutation test, however, returns $p=0.22$
($n_{\mathrm{perm}}=100$), far from the 0.05 significance threshold.
For porosity($\alpha$), random forest reaches $+19.2\%$ RMSE
improvement over the naive-mean baseline (25.380 vs. 31.399,
porosity in percentage points) -- clearing the RMSE bar -- with
permutation $p=0.36$. One of the five porosity candidates, linear
regression, underperforms the baseline outright ($-37.3\%$); the
remaining three learned models cluster closely together and all
clear the RMSE bar (Gaussian process $+18.1\%$, gradient-boosted
trees $+17.9\%$, random forest $+19.2\%$, the best performer),
underscoring that the RMSE improvement random forest achieves is a
genuine, model-class-independent effect and not an artifact of a
single weak comparator.
Refractive index ($+16.2\%$, $p=0.54$) shows the same pattern. Band
gap $E_g(\alpha)$ is the sole observable for which no learned model
even clears the RMSE bar (best learned model $-9.3\%$, i.e.\ worse
than baseline; the naive-mean baseline itself remains the best
``model''), so its permutation test is not applicable.

\begin{table}[htbp]
  \centering
  \caption{Best-model \lomo{}-CV comparison against the pre-registered
    bar, all four target observables, post-comprehensive-audit corpus
    snapshot. RMSE improvement is relative to the tangent-rule baseline
    (\beto{}) or the \lomo{} train-mean baseline (porosity, $n$,
    $E_g$). Permutation $p$ from $n_{\mathrm{perm}}=100$ material-block
    shuffles.}
  \label{tab:pooled_all}
  \begin{tabular}{lrrlrrc}
    \toprule
    Observable & Rows & Materials & Best model & RMSE $\Delta$ & Perm. $p$ & Bar cleared? \\
    \midrule
    $\beta(\alpha)$          & 178 & 37 & gradient-boosted trees & $+49.2\%$ & 0.22 & No \\
    Porosity $\phi(\alpha)$  & 91  & 19 & random forest          & $+19.2\%$ & 0.36 & No \\
    $n(\alpha)$              & 78  & 13 & Gaussian process       & $+16.2\%$ & 0.54 & No \\
    $E_g(\alpha)$            & 67  & 10 & naive-mean baseline    & $+0.0\%$  & n/a  & No \\
    \bottomrule
  \end{tabular}
\end{table}

\subsubsection{Feature importance for the confirmed \beto{} model}
\label{sssec:lofo}

A leave-one-feature-out (LOFO) analysis on the primary Gaussian process
candidate at an earlier corpus snapshot ($n=73$ rows/23 materials,
full-model \lomo{} RMSE $=11.593^{\circ}$) shows the deposition angle
$\alpha$ itself dominates, as expected (RMSE rises to
$19.363^{\circ}$ when removed). Among material descriptors, melting
point (RMSE $13.980^{\circ}$ without it) and atomic mass
($12.585^{\circ}$) carry essentially all of the genuine cross-material
signal beyond $\alpha$; bulk density and crystal structure do not
merely fail to help -- removing either one \emph{improves} \lomo{}
RMSE slightly (to $11.460^{\circ}$ and $10.793^{\circ}$
respectively), indicating these two features contribute noise rather
than signal to the fitted GP at this corpus size
(Table~\ref{tab:lofo}). This descriptive analysis is not itself
subject to the pre-registered significance bar (which applies only to
the model-vs-baseline comparison), but it is directly consistent with
the Bayesian hierarchical model's independent identification of
atomic mass as the one descriptor with a credibly non-zero effect
(Section~\ref{ssec:bayes}).

\begin{table}[htbp]
  \centering
  \caption{Leave-one-feature-out importance for the \beto{} Gaussian
    process model ($n=73$ rows, 23 materials). A larger RMSE increase
    when a feature is removed indicates more predictive signal
    currently attributed to that feature.}
  \label{tab:lofo}
  \begin{tabular}{lrr}
    \toprule
    Feature removed & \lomo{} RMSE without it ($^\circ$) & $\Delta$RMSE \\
    \midrule
    $\alpha$ (deposition angle) & 19.363 & $+7.769$ \\
    Melting point & 13.980 & $+2.387$ \\
    Atomic mass & 12.585 & $+0.991$ \\
    Deposition method & 11.819 & $+0.225$ \\
    \emph{(full model)} & 11.593 & --- \\
    Bulk density & 11.460 & $-0.134$ \\
    Crystal structure & 10.793 & $-0.801$ \\
    \bottomrule
  \end{tabular}
\end{table}

\subsection{The corpus-growth campaign: RMSE improves, significance does not}
\label{ssec:corpus_growth}

Following the initial cleared result at $n=73$ rows/23 materials
($p=0.020$), further, real corpus growth caused the permutation
$p$-value to rise, first to 0.13 at $n=93/25$ materials, confirmed at
a second random seed ($p=0.09$). Five independent corpus-growth
avenues were then pursued deliberately, rather than treating this as
a one-off setback: a targeted literature-backlog absorption pass; a
web-escalation search for entirely new source materials (finding two,
SiO$_2$ and HfO$_2$, with genuine open-access angle-series data, and
correctly excluding two paywalled/snippet-only leads, Ta$_2$O$_5$ and
TiN, rather than fabricating values from partial information); a
feature-engineering ablation adding Cordero covalent atomic radius and
homologous temperature $T/T_{\mathrm{melt}}$; a deep re-digitization of
nine already-absorbed source papers under a stricter protocol of
visually analyzing every figure directly rather than relying on text
extraction (which recovered 43 new data points, most notably growing
ZnO's series from 1 to 22 points); and a whole-drive sweep for
misplaced literature PDFs (which found one further relevant paper,
reporting only two discrete angles, too sparse to contribute \lomo{}
rows). Table~\ref{tab:corpus_growth} reports the resulting best
figures at $n=132$ rows/27 materials.

\begin{table}[htbp]
  \centering
  \caption{Corpus-growth campaign result, all four observables, at
    the $n=132$ rows/27-materials $\beta(\alpha)$ snapshot (the single
    most effective lever tried was the deep figure-level
    re-digitization pass).}
  \label{tab:corpus_growth}
  \begin{tabular}{lrrlrrc}
    \toprule
    Observable & Rows & Materials & Best model & RMSE $\Delta$ & Perm. $p$ & Bar cleared? \\
    \midrule
    $\beta(\alpha)$          & 132 & 27 & gradient-boosted trees & $+46.4\%$ & 0.09 & No \\
    Porosity $\phi(\alpha)$  & 64  & 13 & random forest          & $+28.5\%$ & 0.40 & No \\
    $n(\alpha)$              & 37  & 7  & linear regression      & $+13.1\%$ & 0.70 & No (fails RMSE too) \\
    $E_g(\alpha)$            & 49  & 6  & naive-mean baseline    & ---       & n/a  & No \\
    \bottomrule
  \end{tabular}
\end{table}

The feature-engineering ablation moved neither RMSE improvement nor
permutation $p$ for either $\beta$ or porosity, ruling out ``wrong
feature set'' as an explanation for the shortfall. $\beta(\alpha)$'s
permutation $p$ fell monotonically with every genuine data-density
increase across this round of the investigation
($0.13 \rightarrow 0.09$), a directional trend more consistent with a
real, currently under-sampled effect than with pure noise -- offered
here as a plausibility argument, not as evidence the effect is
proven; the honest reportable result at this snapshot remained NOT
CLEARED.

\subsection{Does material-class heterogeneity explain the pooled failure?}
\label{ssec:stratification}

By a later corpus snapshot -- $n=37$ materials for the pooled
$\beta(\alpha)$ fit, gradient-boosted trees $+49.2\%$, $p=0.22$
(the snapshot in Table~\ref{tab:pooled_all}) -- a natural hypothesis
was tested directly: perhaps pooling elemental metals, oxides/
fluorides, and chalcogenides into a single fit averages away real,
class-specific physics that a per-class fit would recover.
Table~\ref{tab:stratified} reports the pre-registered per-class
result. The hypothesis is \textbf{not confirmed}: no class clears
$p<0.05$. Chalcogenides/sulfides reach $+23.1\%$ RMSE improvement at
$p=0.86$; oxides/fluorides $+41.8\%$ at $p=0.55$; elemental metals --
the largest and chemically most homogeneous subgroup -- $+44.6\%$ at
$p=0.20$, the closest of the three to significance but still far above
the bar. Rock-salt nitrides (4 materials, 18 rows) fell below the
pre-registered five-material minimum and were not tested.

\begin{table}[htbp]
  \centering
  \caption{Material-class-stratified \beto{} results, pre-registered
    grouping by crystal-structure/bonding chemistry, identical
    four-model/\lomo{}/permutation protocol as the pooled fit.
    Pooled reference row included for comparison.}
  \label{tab:stratified}
  \begin{tabular}{lrrlrrc}
    \toprule
    Class & Materials & Rows & Best model & RMSE $\Delta$ & Perm. $p$ & Bar cleared? \\
    \midrule
    Rock-salt nitrides       & 4  & 18 & --- & --- & --- & Not tested ($<5$ materials) \\
    Chalcogenides/sulfides   & 6  & 38 & random forest          & $+23.1\%$ & 0.86 & No \\
    Oxides/fluorides         & 11 & 67 & random forest          & $+41.8\%$ & 0.55 & No \\
    Elemental metals         & 14 & 53 & random forest          & $+44.6\%$ & 0.20 & No \\
    \midrule
    \emph{Pooled (all classes)} & 37 & 178\footnotemark & gradient-boosted trees & $+49.2\%$ & 0.22 & No \\
    \bottomrule
  \end{tabular}
\end{table}
\footnotetext{Row count for the pooled reference row reflects the
  Table~\ref{tab:pooled_all} snapshot; the exact row count at the
  $n=37$-material stratification snapshot differs marginally as
  literature absorption continued between the two runs, per-class row
  counts above are the values actually used for the stratified fits.}

RMSE improvement is large in every class (23--45\%), mirroring the
pooled result: the models find \emph{some} systematic structure in
every partition tried, but that structure does not survive random
re-partitioning of materials into cross-validation folds, in any
grouping. This is not evidence against class-dependent physics -- it
means the current corpus cannot distinguish that hypothesis from
ordinary small-sample noise, since every per-class subgroup (14, 11,
or 6 materials) has fewer materials than the already-underpowered
pooled set (37), trading away exactly the degrees of freedom the
permutation test needs without a compensating gain in signal purity.
Elemental metals' comparatively low $p=0.20$ was flagged as the most
promising lead and pursued further.

\subsection{A targeted corpus-growth push within elemental metals: also null}
\label{ssec:targeted_growth}

Two further, targeted literature-absorption interventions added real,
sourced data specifically to the elemental-metals subgroup, run
separately to isolate which kind of densification helps.
\textbf{Cross-source densification} (adding 1--2 points from
different labs/papers to each of the sparsest metals: W
$1\rightarrow 8$ points across two independent sources, Mo
$1\rightarrow 8$, Au $1\rightarrow 4$) grew the subgroup to 72 rows at
the same 14 materials. RMSE improvement continued to rise, to
$+52.7\%$, but $p$ \emph{regressed}, from 0.20 to 0.30.
\textbf{Same-source densification} (digitizing Figure 11 of
\citet{besnard2025}, a single consistent experimental campaign reporting nine
transition metals on one angle series, adding three genuinely new
materials -- V, Hf, Ta -- to the pre-registered elemental-metals class
under the same fcc/bcc/hcp criterion decided before seeing any result)
grew the subgroup to 17 materials, 88 rows, and gave $p=0.24$ --
better than the cross-source attempt but still worse than the
untouched 14-material baseline, and far from 0.05.
Table~\ref{tab:targeted_growth} summarizes.

\begin{table}[htbp]
  \centering
  \caption{Targeted corpus-growth interventions within the
    elemental-metals subgroup. Neither intervention moves $p$
    reliably toward significance; the range 0.20--0.36 across every
    variant is consistent with ordinary sampling noise around a true
    $p$ that is not small at this sample size.}
  \label{tab:targeted_growth}
  \begin{tabular}{lrrrr}
    \toprule
    Variant & Materials & Rows & RMSE $\Delta$ (best model) & Perm. $p$ \\
    \midrule
    Baseline (Section~\ref{ssec:stratification}) & 14 & 53 & $+44.6\%$ & 0.20 \\
    Cross-source densification (W/Mo/Au)          & 14 & 72 & $+52.7\%$ & 0.30 \\
    Same-source densification (+V, Hf, Ta)        & 17 & 88 & $+54.4\%$ & 0.24 \\
    \bottomrule
  \end{tabular}
\end{table}

Across every variant tried, $p$ ranged 0.20--0.36, with no lever --
more points on existing materials, new materials from a scattered set
of sources, or new materials from one internally consistent source --
moving it reliably toward significance. Recovering significance in
this subgroup would need roughly double the current material count
(30+ elemental metals); a subsequent targeted search for 19 further
elemental metals, mostly platinum-group and rare-earth species,
returned usable new data for only one (Sb) via open web search -- a
genuine access limitation of the search tools available for this
project, not proof the underlying literature does not exist.

\subsection{Reformulating the task: ranking instead of regression}
\label{ssec:ranking}

A separate, orthogonal hypothesis is that the pre-registered
exact-value criterion tests the wrong task: predicting a material's
relative rank against another material at a comparable angle is a
scientifically useful and potentially easier target than predicting
its exact value. Table~\ref{tab:ranking} reports the pre-registered
pairwise-ranking result. For $\beta(\alpha)$ ($n=213$ rows, 40
materials, 10{,}394 directed pairs within a $\pm 5^{\circ}$ angle
window), the best model (a gradient-boosted ranker using the full
descriptor set) reaches a pooled \lomo{} pairwise accuracy of 0.746
(Kendall's $\tau=0.492$ \citep{kendall1938rank}) -- comfortably clearing the 0.60 accuracy
threshold and a clear improvement over both the zero-parameter
tangent-rule baseline and an angle-only logistic baseline (both 0.544
accuracy, $\tau=0.088$). The permutation test, however, gives
$p=0.27$: the significance half of the criterion fails by nearly the
same margin as it does for exact-value regression. For porosity($\alpha$)
($n=92$ rows, 20 materials, 2{,}408 directed pairs), the best model (a
linear ranker) reaches 0.725 accuracy ($\tau=0.450$, versus a
0.602/$\tau=0.204$ baseline) but again fails significance, $p=0.28$.

\begin{table}[htbp]
  \centering
  \caption{Pairwise-ranking reformulation results (pre-registered bar:
    pooled \lomo{} pairwise accuracy $\geq 0.60$ and permutation
    $p<0.05$; $n_{\mathrm{perm}}=100$).}
  \label{tab:ranking}
  \begin{tabular}{lrrlrrc}
    \toprule
    Observable & Rows & Materials & Best model & Accuracy & Perm. $p$ & Bar cleared? \\
    \midrule
    $\beta(\alpha)$         & 213 & 40 & gradient-boosted ranker & 0.746 & 0.27 & No \\
    Porosity $\phi(\alpha)$ & 92  & 20 & linear ranker           & 0.725 & 0.28 & No \\
    \bottomrule
  \end{tabular}
\end{table}

This rules out ``wrong task formulation'' as the explanation for the
pooled negative result: reformulating from exact-value prediction to
relative ranking recovers strong, real-looking accuracy gains over the
physics baseline in both observables, yet the same permutation test
that rejects the regression models also rejects the ranking models, by
almost the same margin. The bottleneck is not what the model is asked
to predict; it is the number of independent material-level samples
available to test whether the predicted structure is more than a
fold-specific artifact.

\subsection{The Bayesian hierarchical partial-pooling model: a genuine complementary positive finding}
\label{ssec:bayes}

As a final, deliberately different line of inquiry, the pre-registered
Bayesian hierarchical model (Section~\ref{ssec:bayes_method}) was fit
to the $\beta(\alpha)$ corpus ($n=213$ rows, 40 materials). All MCMC
diagnostics are clean: maximum $\hat{R} = 1.0000$, minimum bulk
effective sample size $= 1342$, across all seven population-level
parameters.

Of the six descriptor coefficients (three each for the intercept and
slope layers), one -- the effect of atomic mass on a material's mean
offset from the tangent rule, $\mathrm{mean}=-3.639$, 94\% HDI
$=[-7.240,-0.344]$ -- excludes zero. The other five (melting point and
bulk density on both intercept and slope, and mass on slope) have HDIs
that include zero, though several are one-sided enough to be
suggestive rather than flatly null (Table~\ref{tab:bayes_coef}).

\begin{table}[htbp]
  \centering
  \caption{Bayesian hierarchical model: population-level descriptor
    coefficients (standardized descriptors), 94\% highest-density
    intervals (HDI). Only atomic mass's intercept-layer coefficient
    excludes zero.}
  \label{tab:bayes_coef}
  \begin{tabular}{llrrc}
    \toprule
    Layer & Descriptor & Mean & 94\% HDI & Excludes 0? \\
    \midrule
    Intercept ($a_i$) & Atomic mass    & $-3.639$ & $[-7.240,-0.344]$ & \textbf{Yes} \\
    Intercept ($a_i$) & Melting point  & $-3.288$ & $[-6.588, 0.151]$ & No \\
    Intercept ($a_i$) & Bulk density   & $ 1.121$ & $[-2.424, 4.884]$ & No \\
    Slope ($b_i$)      & Atomic mass    & $-1.554$ & $[-4.916, 1.834]$ & No \\
    Slope ($b_i$)      & Melting point  & $-2.552$ & $[-5.943, 0.475]$ & No \\
    Slope ($b_i$)      & Bulk density   & $ 2.214$ & $[-1.070, 5.528]$ & No \\
    \bottomrule
  \end{tabular}
\end{table}

The primary deliverable, the Bayesian $R^2$, is more informative than
the per-coefficient test: the posteriors are not concentrated near
zero. The intercept-layer $R^2_a$ has posterior median 0.28 (94\% HDI
$[0.045,0.512]$; $P(R^2_a>0.1)=92\%$, $P(R^2_a>0.2)=72\%$); the
slope-layer $R^2_b$ has median 0.20 (94\% HDI $[0.005,0.430]$;
$P(R^2_b>0.1)=80\%$, $P(R^2_b>0.2)=50\%$). Read plainly: it is quite
likely (80--92\% posterior probability) that these three descriptors
explain a non-trivial share of material-to-material variation in
$\beta(\alpha)$, with better-than-even odds the share exceeds a fifth
-- but the credible intervals are wide and do not exclude a near-zero
effect at one end, and substantial residual per-material idiosyncrasy
remains unexplained either way (residual scales $\sigma_a=8.41^{\circ}$,
$\sigma_b=7.03^{\circ}$, both large relative to the mean coefficients).

This does not contradict the negative frequentist verdict of
Sections~\ref{ssec:pooled}--\ref{ssec:ranking}, and should not be read
as doing so -- it answers a different question. A real-but-modest
effect, imprecisely estimated (median $R^2$ around 0.2--0.3 but a wide
HDI reaching close to zero), is exactly the kind of signal a
leave-one-material-out permutation test at $n=37$--40 materials is
underpowered to detect at $p<0.05$ even when the effect is genuine.
Two caveats apply. First, crystal structure was excluded from the
group-level regression for identifiability reasons (19 levels across
43 materials); if it carries real information, as is physically
plausible for a shadowing/anisotropy effect, the true
descriptor-explainable variance could be higher than reported here.
Second, the Bayesian $R^2$ reported is in-sample, not
cross-validated -- it measures how much structure the fitted model
captures, not how well it would predict a new material's curve.

\subsection{Summary of results}
\label{ssec:results_summary}

Across five independent lines of attack -- the pooled model at a
further-grown corpus, material-class stratification, a targeted
corpus-growth push confined to the most promising subgroup, a
ranking/ordinal task reformulation, and a Bayesian hierarchical
framework -- no frequentist formulation clears the pre-registered
significance bar. RMSE (or pairwise-ranking accuracy) improvement
over the physics baseline is real and large in every configuration
tested, from $+23\%$ to $+54\%$ depending on subgroup, and 0.72--0.75
pairwise accuracy against a 0.54--0.60 baseline for the ranking
reformulation, so the learned models consistently find genuine,
reproducible structure in the data rather than overfitting to noise
that happens to look good in-sample. The barrier throughout is
specifically statistical power, not model quality or task framing:
every reformulation tried reproduces the same pattern of strong effect
size and weak significance, and the one framework built to be
sensitive to a real-but-modest effect under small samples (the
Bayesian hierarchical model) finds exactly that.

% ============================================================
\section{Discussion}
\label{sec:discussion}

\subsection{Why this is a genuine finding, not a failed project}
\label{ssec:disc_genuine}

A reader could reasonably ask whether a paper reporting ``no
configuration clears a pre-registered significance bar'' is worth
publishing. We argue it is, for three reasons specific to this
investigation's design, not as a general defense of negative results.

First, the diagnostic completeness is itself informative. Four
distinct, plausible explanations for the pooled model's shortfall were
tested directly and independently ruled out: wrong feature set (the
feature-engineering ablation of Section~\ref{ssec:corpus_growth}
changed neither RMSE improvement nor $p$), wrong extraction method
(the deep figure-re-digitization pass of the same section improved $p$
from 0.13 to 0.09 but did not close the gap), class heterogeneity
(Section~\ref{ssec:stratification}, no subgroup clears the bar despite
each being chemically more homogeneous than the pooled set), and wrong
task formulation (Section~\ref{ssec:ranking}, the ranking
reformulation reproduces the identical strong-effect/weak-significance
pattern). Ruling out four specific, testable alternative explanations
and converging on a fifth (sample size) that is directly measurable
and directly actionable is a stronger scientific contribution than an
unexplained single negative result would be.

Second, the improving trend is consistent with a real, recoverable
effect rather than pure noise. $\beta(\alpha)$'s permutation $p$ fell
monotonically with every genuine data-density increase in the
corpus-growth campaign of Section~\ref{ssec:corpus_growth}
($0.13\rightarrow 0.09$), and the material-class-stratified subgroup
closest to significance (elemental metals) is also the chemically most
homogeneous one -- directionally consistent with ``more physically
similar materials, and more data on the same materials, moves $p$ the
right way,'' even though neither trend alone reaches significance at
current sample sizes. We present this as a plausibility argument for
future corpus growth, not as proof the effect is real; the honest
frequentist verdict at every snapshot tested here remains NOT CLEARED.

Third, and most directly, the Bayesian hierarchical analysis
(Section~\ref{ssec:bayes}) demonstrates \emph{why} a permutation-test
failure at $n=14$--40 materials does not mean there is no effect. The
same corpus that fails the material-level permutation test at
$p=0.20$--0.36 yields Bayesian $R^2$ posteriors with 80--92\%
probability of explaining more than 10\% of between-material variance
in $\beta(\alpha)$, concentrated specifically on atomic mass. These
two results are not in tension: a permutation test at this sample size
is a coarse, conservative instrument, and a real-but-modest effect of
exactly the size the Bayesian model estimates is precisely what it is
underpowered to detect. Reporting only the frequentist NOT-CLEARED
verdict, without the Bayesian analysis that characterizes \emph{how}
partial the evidence is, would understate what this corpus actually
supports; reporting only the Bayesian result, without the frequentist
verdict, would overstate it. Presenting both together is, we argue,
the honest way to report this class of result.

\subsection{Sample-size, not extraction quality, is the binding constraint}
\label{ssec:disc_samplesize}

The single most consistent pattern across every section of
Section~\ref{sec:results} is that RMSE improvement (or pairwise-ranking
accuracy) is never the limiting factor -- it clears its threshold in
every configuration tested except $E_g(\alpha)$, which independently
fails for a different, likely corpus-specific reason (only 6--10
materials, the smallest of the four observables). The limiting factor
is always the number of \emph{independent materials} available to the
permutation test, which sets the resolution of its null distribution.
This has a direct, testable implication: the levers that do not work
(more rows on existing materials from scattered sources, different
task formulations, different material-class partitions) all leave the
material count roughly unchanged or reduce it; the one lever that
measurably improved $p$ within this investigation (same-source
figure-level re-digitization, which both added rows and, when it
introduced genuinely new materials like the Besnard densification, did
so cleanly) is the one that increases the corpus's material diversity
without simultaneously increasing between-source measurement variance.
This is consistent with the cross-source-vs-same-source contrast
observed directly in Section~\ref{ssec:targeted_growth}: adding W, Mo,
and Au points from scattered, independent labs increased row count but
\emph{worsened} $p$ (0.20$\rightarrow$0.30), plausibly because
different labs' deposition conditions (rate, substrate, base pressure)
introduce real between-source variance that a single-descriptor model
cannot separate from genuine material effect, whereas the
single-campaign Besnard extension added new materials with internally
consistent measurement conditions and moved $p$ in the right direction
relative to the cross-source attempt (though still short of the
14-material baseline).

\subsection{What would resolve the open question}
\label{ssec:disc_resolve}

The material-class-stratification result identifies the most
promising, concrete target for future work: elemental metals reach
$p=0.20$ at $n=14$ materials, the closest of any configuration tested
to the 0.05 bar, and are also the chemically most homogeneous and
best-characterized class (complete, high-confidence descriptor
coverage for all four descriptors used). Recovering significance in
this subgroup plausibly requires roughly double the current material
count -- on the order of 30+ elemental metals with genuine
angle-resolved $\beta(\alpha)$ data. A targeted search for 19 further
elemental metals, concentrated on platinum-group and rare-earth
species not yet in the corpus, returned usable data for only one (Sb)
via the open web search tools available to this project; this is a
practical ceiling of accessible open literature search, not evidence
that the underlying data does not exist in the broader, harder-to-reach
literature (older or non-open-access journals, conference proceedings,
and materials-science theses not indexed by general web search).
Access to a professional literature database with more complete
coverage of these specific material families (Scopus, Web of Science,
or equivalent) is therefore the single most direct, well-motivated
path to resolving whether the elemental-metals signal is real. We
flag this explicitly as the concrete deliverable this paper offers to
the next phase of this research program, rather than leaving ``more
data would help'' as an unspecific closing remark.

\subsection{Limitations}
\label{ssec:disc_limitations}

Several limitations apply beyond the sample-size discussion above.
The corpus is drawn from heterogeneous deposition conditions (rate,
base pressure, substrate, chamber geometry) that are not uniformly
reported across source papers and are only partially captured by the
deposition-method categorical descriptor; residual between-source
variance not attributable to material identity is folded into what
the permutation test treats as noise, which could mask a real effect
or, less plausibly, contribute to the apparent RMSE improvements
through incidental correlation with material identity. The Bayesian
hierarchical model's $R^2$ is in-sample, not predictive, and excludes
crystal structure for identifiability reasons; both caveats mean its
positive finding should be read as a characterization of the corpus's
explainable structure, not a validated predictive claim. Finally, all
results here concern angle-resolved morphological observables
($\beta$, porosity) under a single physics baseline (the tangent
rule); the two companion observables tested (refractive index, band
gap) are optical rather than morphological and were evaluated
primarily to check whether the negative pattern was specific to
$\beta$/porosity or general across GLAD-relevant observables -- it
proved general, which we read as additional, independent evidence
that the underlying constraint is sample size rather than an
observable-specific modeling failure.

% ============================================================
\section{Conclusion}
\label{sec:conclusion}

We conducted a pre-registered, leakage-controlled evaluation of
whether material descriptors (atomic mass, melting point, bulk
density, crystal structure, deposition method) can predict a new
material's column-tilt angle $\beta(\alpha)$ or porosity $\phi(\alpha)$
under glancing-angle deposition, using leave-one-material-out
cross-validation and a material-level permutation test as the
significance instrument. Every learned model comfortably clears the
pre-registered RMSE-improvement bar for $\beta(\alpha)$ (up to
$+54.4\%$ depending on corpus snapshot) and usually for porosity (up
to $+28.5\%$), but no configuration tested -- the pooled fit across up
to 40 materials, three pre-registered material-class subgroups, two
independent targeted corpus-growth interventions, or a pairwise
ranking reformulation of the task -- ever clears the permutation-test
significance bar. Four plausible alternative explanations (feature
choice, extraction quality, class heterogeneity, task formulation)
were tested directly and ruled out; the specific, remaining, and
well-quantified bottleneck is material-level sample size. A
complementary Bayesian hierarchical partial-pooling model, fit
independently and asking a different statistical question, finds a
genuine, credible non-zero effect of atomic mass on the material-level
offset from the physics baseline and places 80--92\% posterior
probability on descriptors explaining more than 10\% of
between-material variance -- a result fully consistent with, and
explanatory of, the frequentist non-significance verdict rather than
contradicting it.

We report this as a rigorous, multi-angle-tested negative-to-mixed
result: a well-instrumented account of exactly where the
statistical-power boundary for cross-material GLAD morphology
prediction currently sits, together with the specific, actionable
corpus-growth target (professional literature-database access to
30+ elemental metals, particularly platinum-group and rare-earth
species) most likely to move it. We offer this both as a direct
contribution -- ruling out four alternative explanations for a
persistent gap between effect size and statistical significance in
small-corpus materials-informatics problems is itself useful to
readers attempting similar cross-material generalization studies --
and as a resource for future corpus-growth efforts specifically in
this sub-field, so that the next research effort in this direction can
begin from a documented, tested starting point rather than repeating
the same exploratory search.

% ============================================================
\section*{Data and Code Availability}
All feature tables, model-comparison scripts, and pre-registration
documents referenced in this paper are maintained in the project
repository (\texttt{02\_PINN\_NIF/P3\_ML/}); the corpus is drawn from
peer-reviewed literature and PhD theses, cited where individual source
papers contributed the majority of a material's data
(Section~\ref{sec:intro}), with the full per-row provenance recorded
in the underlying feature-table generation pipeline.

% ============================================================
\section*{CRediT Author Statement}
Oussama Ben Haj Ali: conceptualization, methodology, software,
formal analysis, investigation, data curation, writing -- original
draft, writing -- review and editing, visualization. Ferid Chaffar
Akkari, Bruno Gallas, Nicolas Martin: resources, supervision, writing
-- review and editing.

\bibliographystyle{aipnum4-1}
\bibliography{P3_references}

\end{document}
